\documentclass[11pt]{article}

\usepackage{amsmath,amssymb,amsthm}

\newtheorem{theorem}{Theorem}
\newtheorem{lemma}{Lemma}
\newtheorem{corollary}{Corollary}
\newtheorem{remark}{Remark}

\DeclareMathOperator{\ch}{ch}

\begin{document}
\title{Submission C solution to Problem 6}
\date{}
\maketitle


\begin{remark}
Throughout, the phrase ``a single vertex'' with \(w(v)<d_T(v)\) is interpreted as saying that \(v\) is the unique such vertex.
\end{remark}

\begin{lemma}[Norm-one elements are irreducible]
Let \(L\) be a positive definite integral lattice. If \(x\in L\) satisfies \(x^2=1\), then \(x\) is irreducible.
\end{lemma}

\begin{proof}
Suppose \(x=a+b\), where \(a,b\in L\) are nonzero and \(a\cdot b\ge 0\). Since \(L\) is positive definite and integral, \(a^2,b^2\in \mathbb Z_{>0}\). Hence
\[
1=x^2=a^2+b^2+2a\cdot b\ge 1+1=2,
\]
a contradiction.
\end{proof}

Let \(C\) be a rooted weighted tree with root \(\rho\). Orient every edge away from \(\rho\), and write \(\ch_C(x)\) for the number of children of a vertex \(x\in C\). We call \(C\) admissible if its form is positive definite, every vertex weight is at least \(2\), and
\[
w(x)\ge \ch_C(x)+1
\qquad \text{for every } x\in C.
\]
For \(x=\sum n_y y\in L(C)\), write \(x_y=n_y\) for the coefficient of the vertex \(y\).

For a positive definite rooted tree \(C\), define its capacity by
\[
\gamma(C)=\bigl(Q_C^{-1}\bigr)_{\rho\rho},
\]
where \(Q_C\) is the Gram matrix of \(C\) in the vertex basis.

\begin{lemma}[Capacities of admissible rooted trees]
Let \(C\) be an admissible rooted tree with root \(\rho\). Then
\[
0<\gamma(C)<1.
\]
Moreover, if the root has children \(\rho_1,\dots,\rho_t\), and \(C_i\) denotes the rooted subtree with root \(\rho_i\), then
\[
\gamma(C)=\frac{1}{w(\rho)-\sum_{i=1}^t \gamma(C_i)}.
\]
\end{lemma}

\begin{proof}
Order the vertices so that \(\rho\) comes first and the remaining vertices are grouped by the rooted subtrees \(C_1,\dots,C_t\). Then
\[
Q_C=
\begin{pmatrix}
w(\rho) & c^T\\
c & Q_0
\end{pmatrix},
\]
where \(Q_0=\bigoplus_i Q_{C_i}\), and \(c\) has entry \(-1\) in the coordinate corresponding to each child root \(\rho_i\) and \(0\) elsewhere. By the Schur complement formula,
\[
\bigl(Q_C^{-1}\bigr)_{\rho\rho}
=
\frac{1}{w(\rho)-c^TQ_0^{-1}c}
=
\frac{1}{w(\rho)-\sum_i \gamma(C_i)}.
\]

It remains to show that \(\gamma(C)<1\). This follows by induction on the number of vertices. If \(C\) has one vertex, then \(w(\rho)\ge 2\), so \(\gamma(C)=1/w(\rho)<1\). Otherwise, by induction, \(\gamma(C_i)<1\) for every child subtree. If the root has \(t\) children, then admissibility gives \(w(\rho)\ge t+1\). Hence
\[
w(\rho)-\sum_i \gamma(C_i)>w(\rho)-t\ge 1,
\]
so \(\gamma(C)<1\).
\end{proof}

\begin{lemma}[A rooted estimate]
Let \(C\) be an admissible rooted tree with root \(\rho\) and capacity \(\gamma=\gamma(C)\). Then, for every \(x\in L(C)\) and every \(k\in \mathbb Z\),
\[
x^2-(2k+1)x_\rho+\gamma k(k+1)\ge 0.
\]
\end{lemma}

\begin{proof}
We prove the claim by induction on \(|C|\).

If \(C\) has one vertex, say \(x=m\rho\), then \(\gamma=1/w(\rho)\). Let \(W=w(\rho)\). Multiplying the desired inequality by \(W\), we get
\[
W\left(Wm^2-(2k+1)m+\frac{k(k+1)}{W}\right)
=
(Wm-k)(Wm-k-1).
\]
Since \(Wm-k\in\mathbb Z\), the product of the two consecutive integers \(Wm-k\) and \(Wm-k-1\) is nonnegative. This proves the base case.

Now assume \(|C|>1\). Let the root be \(\rho\), let its children be \(\rho_1,\dots,\rho_t\), and let \(C_i\) be the rooted subtree with root \(\rho_i\). Write
\[
x=a\rho+\sum_i x_i,
\qquad x_i\in L(C_i),
\]
and let \(s_i=(x_i)_{\rho_i}\). Then
\[
x^2=w(\rho)a^2-2a\sum_i s_i+\sum_i x_i^2.
\]
Let \(\gamma_i=\gamma(C_i)\). By the induction hypothesis applied to \(C_i\) with \(k=a\), we have
\[
x_i^2-(2a+1)s_i+\gamma_i a(a+1)\ge 0,
\]
so
\[
x_i^2-2as_i\ge s_i-\gamma_i a(a+1).
\]
Similarly, applying the induction hypothesis with \(k=a-1\), we obtain
\[
x_i^2-(2a-1)s_i+\gamma_i a(a-1)\ge 0,
\]
so
\[
x_i^2-2as_i\ge -s_i-\gamma_i a(a-1).
\]
Taking the maximum of these two lower bounds gives
\[
x_i^2-2as_i
\ge
-\gamma_i a^2+\bigl|s_i-\gamma_i a\bigr|.
\]
Therefore, setting
\[
D=\sum_i \bigl|s_i-\gamma_i a\bigr|,
\]
we get
\begin{align*}
x^2-(2k+1)a+\gamma k(k+1)
&\ge
\left(w(\rho)-\sum_i\gamma_i\right)a^2-(2k+1)a+\gamma k(k+1)+D.
\end{align*}
By the capacity formula,
\[
\gamma^{-1}=w(\rho)-\sum_i\gamma_i.
\]
Thus, if \(\tau=a/\gamma\), the preceding lower bound becomes
\[
\gamma(\tau-k)(\tau-k-1)+D.
\]

If \(\tau\notin(k,k+1)\), then \((\tau-k)(\tau-k-1)\ge 0\), and the result follows. Suppose instead that \(k<\tau<k+1\). Put
\[
\alpha=\tau-k,
\qquad
\beta=k+1-\tau.
\]
Then \(0<\alpha,\beta<1\), \(\alpha+\beta=1\), and
\[
\gamma(\tau-k)(\tau-k-1)=-\gamma\alpha\beta.
\]
Now
\[
\tau
=
\frac{a}{\gamma}
=
\left(w(\rho)-\sum_i\gamma_i\right)a.
\]
Since \(w(\rho)a-\sum_i s_i\in\mathbb Z\), we have
\[
D
=
\sum_i |s_i-\gamma_i a|
\ge
\operatorname{dist}(\tau,\mathbb Z)
=
\min\{\alpha,\beta\}.
\]
Because \(0<\gamma<1\) and \(\alpha\beta\le \min\{\alpha,\beta\}\), it follows that
\[
D\ge \gamma\alpha\beta.
\]
Hence
\[
\gamma(\tau-k)(\tau-k-1)+D\ge 0.
\]
This completes the induction.
\end{proof}

\begin{corollary}
Let \(C\) be an admissible rooted tree with root \(\rho\). If \(0\ne x\in L(C)\), then
\[
x^2-x_\rho>0.
\]
\end{corollary}

\begin{proof}
Let \(\gamma=\gamma(C)\). If \(x_\rho\le 0\), then \(x^2-x_\rho>0\) because \(C\) is positive definite and \(x\ne 0\).

Assume \(x_\rho>0\). Let \(\rho^\#\in L(C)\otimes \mathbb Q\) be the vector satisfying
\[
\rho^\#\cdot y=y_\rho
\qquad
\text{for all } y\in L(C).
\]
Equivalently, the coefficient vector of \(\rho^\#\) is \(Q_C^{-1}e_\rho\). Thus
\[
(\rho^\#)^2=\gamma.
\]
By Cauchy's inequality in the positive definite real vector space \(L(C)\otimes\mathbb R\),
\[
x_\rho^2=(x\cdot \rho^\#)^2\le x^2(\rho^\#)^2=\gamma x^2.
\]
Since \(\gamma<1\), we get
\[
x^2\ge \frac{x_\rho^2}{\gamma}>x_\rho^2\ge x_\rho.
\]
Therefore \(x^2-x_\rho>0\).
\end{proof}

\begin{theorem}
Let \(T\) be a finite weighted tree whose associated form is positive definite. Suppose that there is a unique vertex \(v\) such that
\[
w(v)<d_T(v).
\]
Then \(T\) contains a vertex that is irreducible in \(L(T)\).
\end{theorem}

\begin{proof}
Since the form is positive definite, every vertex has positive square. Thus every vertex weight is a positive integer.

If some vertex \(u\) has \(w(u)=1\), then \(u^2=1\), and \(u\) is irreducible by the first lemma. Hence we may assume from now on that every vertex weight is at least \(2\).

Let \(v\) be the unique vertex with \(w(v)<d_T(v)\). We will prove that \(v\) itself is irreducible.

Let \(C_1,\dots,C_m\) be the connected components of \(T\setminus\{v\}\). For each \(i\), let \(\rho_i\) be the unique vertex of \(C_i\) adjacent to \(v\), and root \(C_i\) at \(\rho_i\). Because \(v\) is the unique bad vertex, every vertex \(x\in C_i\) satisfies
\[
w(x)\ge d_T(x).
\]
With the orientation away from \(\rho_i\), each \(x\in C_i\) has exactly one edge pointing toward \(v\) or toward its parent, so
\[
d_T(x)=\ch_{C_i}(x)+1.
\]
Therefore
\[
w(x)\ge \ch_{C_i}(x)+1.
\]
Since all weights are at least \(2\), each \(C_i\) is admissible. Let
\[
\gamma_i=\gamma(C_i).
\]

By the Schur complement formula applied to the decomposition of \(T\) into \(v\) and the components \(C_i\), positive definiteness of \(T\) implies
\[
w(v)-\sum_i\gamma_i>0.
\]
Write \(W=w(v)\).

Suppose, toward a contradiction, that \(v\) is reducible. Then there exist nonzero \(a,b\in L(T)\) such that
\[
v=a+b,
\qquad
a\cdot b\ge 0.
\]
Put \(z=-b\). Then \(a=v+z\), and
\[
0\le a\cdot b=(v+z)\cdot(-z)=-(z^2+v\cdot z).
\]
Hence
\[
z^2+v\cdot z\le 0.
\]
Also \(z\ne 0\) and \(z\ne -v\), because \(a,b\ne 0\).

Let \(q\) be the coefficient of \(v\) in \(z\). If \(q\le -1\), replace \(z\) by
\[
z'=-v-z.
\]
This replacement corresponds to interchanging \(a\) and \(b\), and
\[
z'\cdot(z'+v)=z\cdot(z+v).
\]
Moreover, the coefficient of \(v\) in \(z'\) is \(-1-q\ge 0\). Thus we may assume that the coefficient of \(v\) in \(z\) is a nonnegative integer \(p\).

Write
\[
z=pv+\sum_i y_i,
\qquad y_i\in L(C_i).
\]
Let \(s_i=(y_i)_{\rho_i}\). Since \(v\cdot y_i=-s_i\), we compute
\begin{align*}
z^2+v\cdot z
&=
\left(Wp^2-2p\sum_i s_i+\sum_i y_i^2\right)
+
\left(Wp-\sum_i s_i\right) \\
&=
Wp(p+1)+\sum_i\left(y_i^2-(2p+1)s_i\right).
\end{align*}

First suppose \(p\ge 1\). Applying the rooted estimate to each \(C_i\) with \(k=p\), we get
\[
y_i^2-(2p+1)s_i\ge -\gamma_i p(p+1).
\]
Therefore
\[
z^2+v\cdot z
\ge
\left(W-\sum_i\gamma_i\right)p(p+1)>0,
\]
contradicting \(z^2+v\cdot z\le 0\).

It remains to consider \(p=0\). Then
\[
z^2+v\cdot z
=
\sum_i (y_i^2-s_i).
\]
By the corollary, each term \(y_i^2-s_i\) is nonnegative, and it is strictly positive whenever \(y_i\ne 0\). Since \(z\ne 0\) and \(p=0\), at least one \(y_i\) is nonzero. Hence
\[
z^2+v\cdot z>0,
\]
again contradicting \(z^2+v\cdot z\le 0\).

Thus \(v\) is irreducible. Consequently, \(T\) contains an irreducible vertex.
\end{proof}

\end{document}